%% ===================================================================
\section{The finite part}\label{sec:finite}
%% ===================================================================

\subsection{Definition}\label{sec:findef}

Every row of Sec.~\ref{sec:xsec} has the form $\int dp^+\,I_{\rm row}(p^+)$, with $p^+$ the plus momentum of the unobserved or virtual gluon. We define its finite part by the exact subtraction
\begin{equation}
\frac{d^3N^{\rm fin}_{\rm row}}{d^3k}\equiv\frac{d^3N_{\rm row}}{d^3k}-\delta Y_T\,\mathfrak T_{\rm row}-\delta Y_P\,\mathfrak P_{\rm row},\qquad \mathfrak T_{\rm row}=\lim_{p^+\to0}p^+I_{\rm row},\quad \mathfrak P_{\rm row}=\lim_{p^+\to\infty}p^+I_{\rm row}.
\label{eq:findef}
\end{equation}
$\mathfrak T_{\rm row}$ and $\mathfrak P_{\rm row}$ depend neither on $\Lambda$ nor on $\vee$, and $d^3N^{\rm fin}_{\rm row}$ has a finite limit for $\Lambda\to0$, $\vee\to\infty$. Since $\vee$ is a physical, finite cutoff (the boundary between the soft window and the valence modes), we keep every function of $k^+/\vee$ exactly. We also keep the functions of $\Lambda/k^+$ that vanish as $\Lambda\to0$, so that every equation below is an exact identity. The three small logarithms that occur are
\begin{equation}
\lV\equiv\log\Big(1-\frac{k^+}{\vee}\Big),\qquad \lL\equiv\log\Big(1-\frac{\Lambda}{k^+}\Big),\qquad \lLb\equiv\log\Big(1+\frac{\Lambda}{k^+}\Big).
\label{eq:smalllogs}
\end{equation}
Summed over the rows, $\mathfrak T$ and $\mathfrak P$ give JIMWLK$\otimes$LO and BFKL$\otimes$LO of Sec.~\ref{sec:LL}. The finite part of the cross section is $\sum_{\rm rows}d^3N^{\rm fin}_{\rm row}$.

\paragraph*{Two practical rules.}
\begin{itemize}
\item[(A)] \emph{Rows in which the $p^+$ integral has been done.} Every longitudinal function $L$ is split as $L=c_T\,\delta Y_T+c_P\,\delta Y_P+L^{\rm fin}$, with $c_T,c_P$ numbers. The splits are collected in Table~\ref{tab:L} of App.~\ref{app:long}. Rapidity logarithms can be hidden in functions that look like logarithms of transverse distances. For example, with $a,b$ squared transverse distances, $\log\frac{(\vee-k^+)b+k^+a}{\Lambda b+k^+a}\to\log\frac{\vee}{k^+}+\log\frac{b}{a}$ for $\Lambda\ll k^+\ll\vee$, so it contains $\delta Y_P$.
\item[(B)] \emph{Rows left as plus-prescription integrals} (G1-V, G2-IV, G3-III, G3-V, G3-VI). The plus prescription removes the $p^+\to0$ end, and the accompanying endpoint logarithm gives $\mathfrak T$. The end $p^+\to\infty$ is still inside the integral, whose upper limit is $\vee-k^+$. Two mechanisms hide a $\delta Y_P$ there. The first is the plus prescription itself: with a phase under the integral,
\begin{equation}
\int_0^Pdp^+\,\frac{e^{-i\beta p^+}-1}{p^+}=-\log(|\beta|P)-\gamma_E-\frac{i\pi}{2}\,{\rm sgn}\,\beta+\Ord\Big(\frac{1}{\beta P}\Big).
\label{eq:phase}
\end{equation}
For $P=\vee-k^+$ and $\beta=\mathbf k\cdot(y'-z)/k^+$ this contains $-\delta Y_P-\lV-\log|\mathbf k\cdot(y'-z)|$. The second mechanism is a term that looks regular at large $p^+$ in the shifted variable $y'=z+\frac{k^+}{k^++p^+}(w'-z)$ but behaves as $1/p^+$ in the original variable $w'$. For these rows we give $\mathfrak P$ explicitly, and the finite part is Eq.~\eqref{eq:findef}.
\end{itemize}

\paragraph*{Principal value.} Several $p^+$ integrands have a simple pole inside the integration range, for example $1/[p^+(x-z)^2-k^+(x-w)^2]$ at $p^+=k^+(x-w)^2/(x-z)^2$. This pole is the energy denominator of the one-gluon state, and we define the integrals by their principal value. All logarithms of the results below are then to be read as $\log|\cdot|$. With a different prescription, terms $\pm i\pi$ times the coefficient of the pole would appear.

\paragraph*{Notation.} The measured gluon sits at $w$ in the amplitude and at $w'$ in its conjugate; the unobserved gluon sits at $z$. We use the shorthands
\begin{equation}
a\equiv(x-w)^2,\quad b\equiv(x-z)^2,\quad a'\equiv(x'-w')^2,\quad b'\equiv(x'-z)^2,\quad Z^2\equiv(z-w)^2,\quad Z'^2\equiv(z-w')^2,
\end{equation}
and $X\equiv x'-w'$, $Y\equiv y-z$ where indicated. Each row is written as the term displayed, plus its complex conjugate where stated. The colour operators are given with the shorthand $U^{bc}\equiv U^{cb}$ of Sec.~\ref{sec:xsec}.

%% -------------------------------------------------------------------
\subsection{Group I}\label{sec:finiteG1}

\paragraph*{G1-I.} This row contains only four-charge terms, each proportional to exactly $\log(\vee/\Lambda)$. It has no finite part, apart from the edge term discussed in Sec.~\ref{sec:edge}.

\paragraph*{G1-II, part 1 (one-charge part of $\mf A^{(3)}$).} The rapidity logarithms of Eq.~\eqref{eq:grho} combine according to Eq.~\eqref{eq:grhologs} into $-2\delta Y_T+\lV$; there is no $\delta Y_P$. The finite part is
\begin{align}
\frac{d^3N^{\rm fin}_{\rm II,1}}{d^3k}={}&\frac{1}{(2\pi)^3}\frac{g^4N_c}{32\pi^4k^+}\int_{w,w',x,x'}e^{-ik\cdot(w'-w)}\,K_{\rm LO}\,\Bigg\{\Big[-\frac2\epsilon-2\gamma_E-\log\frac{(x-w)^2\bar\mu^2}{4}\Big]\Big[\frac{11}{3}+\lV\Big]\notag\\
&\hspace{10em}-\frac{67}{9}+\frac{\pi^2}{3}+\lV^2+\lV\log\frac{\vee}{k^+}\Bigg\}\,\Ord_{\rm LO}+{\rm c.c.},
\label{eq:finII1}
\end{align}
and the subtracted part is the same expression with the curly bracket replaced by $-2\delta Y_T\big[-\frac2\epsilon-2\gamma_E-\log\frac{(x-w)^2\bar\mu^2}{4}\big]$. The term $\lV\log(\vee/k^+)$ is a large logarithm times $\Ord(k^+/\vee)$; it is bounded and belongs to the finite part. Relative to the LO spectrum~\eqref{eq:LOrelab}, Eq.~\eqref{eq:finII1} is $-\frac{\alpha_sN_c}{4\pi}\{\dots\}$ per ordering. The $\frac{11}{3}$ term is the $\beta_0$ term and $-\frac{67}{9}+\frac{\pi^2}{3}$ the CMW constant (Sec.~\ref{sec:rc}).

\paragraph*{G1-II, part 2 (two-charge part of $\mf A^{(3)}$).} In the mixed representation of Eq.~\eqref{eq:grhorho}, with the Fourier variable called $\mathbf q$ to avoid confusion with the measured $\mathbf k$,
\begin{align}
\frac{d^3N^{\rm fin}_{\rm II,2}}{d^3k}={}&-\frac{1}{(2\pi)^3}\frac{g^4f^{abc}}{64\pi^6k^+}\int_{w,w'}e^{-ik\cdot(w'-w)}\int d^2\mathbf p\,d^2\mathbf q\notag\\
&\times\int_{x,y,x'}e^{i\mathbf q\cdot(w-y)}e^{i\mathbf p\cdot(y-x)}\frac{(x'-w')^i}{(x'-w')^2\,\mathbf p^2\mathbf q^2}\notag\\
&\times\Bigg[-2\mathbf q^i\log\frac{\mathbf p^2}{\mathbf q^2+(k^+/\vee)\,\mathbf p^2}-\frac{\mathbf p^2\mathbf q^i+\mathbf q^2\mathbf p^i}{(\mathbf q-\mathbf p)^2}\,\lV\notag\\
&\qquad+\frac{(2\mathbf p^2-2\mathbf q\cdot\mathbf p+\mathbf q^2)\mathbf q^i+\mathbf q^2\mathbf p^i}{(\mathbf q-\mathbf p)^2}\log\frac{(1-k^+/\vee)\,\mathbf q^2+(k^+/\vee)(\mathbf q-\mathbf p)^2}{(\mathbf q-\mathbf p)^2}\Bigg]\notag\\
&\times\Big[U^{ab'}(x')\rho^{b'}(x')-U^{ab'}(w')\rho^{b'}(x')\Big]\notag\\
&\times\Big[U^{ab}(w)\{\rho^b(x),\rho^c(y)\}-U^{bd}(x)U^{ce}(y)\{\rho^d(x),\rho^e(y)\}\Big]+{\rm c.c.}
\label{eq:finII2}
\end{align}
The subtracted part is $\delta Y_T$ times $-\frac{2\mathbf q^2\mathbf p^i+(\mathbf p^2-\mathbf q\cdot\mathbf p)\mathbf q^i}{(\mathbf q-\mathbf p)^2}-2\mathbf q^i$ in the square bracket. The coefficient of $\delta Y_P$ vanishes identically. Taking the six terms of the bracket of Eq.~\eqref{eq:grhorho} in order, they contribute $0$, $+2\mathbf q^i$, $-2\mathbf q^i$, $+\mathbf q^i$ (the fourth and fifth terms together) and $-\mathbf q^i$, which add up to zero. For $\vee\to\infty$ the square bracket becomes $-2\mathbf q^i\log\frac{\mathbf p^2}{\mathbf q^2}+\frac{(2\mathbf p^2-2\mathbf q\cdot\mathbf p+\mathbf q^2)\mathbf q^i+\mathbf q^2\mathbf p^i}{(\mathbf q-\mathbf p)^2}\log\frac{\mathbf q^2}{(\mathbf q-\mathbf p)^2}$.

\paragraph*{G1-III (rescattering, $\bar{\mf B}_3^\dagger\mf A$).} The colour operator, with the prefactor $f^{abd}$, is
\begin{equation}
\Ord_{\rm III}\equiv\Big[U^{ab'}(x')\rho^{b'}(x')-U^{ab'}(w')\rho^{b'}(x')\Big]\Big[U^{de}(x)U^{bc}(z)\rho^e(x)\rho^c(y)-U^{de}(x)U^{bc}(y)\rho^e(x)\rho^c(y)\Big].
\end{equation}
The row contains the three parts of $\mf B_3$ with one valence charge: the instantaneous term, the non-instantaneous terms of both orderings, and the commutator term. With $X=x'-w'$, $Y=y-z$ define the kernels
\begin{align}
&K_a=\frac{X\cdot Y\,(a+b)}{X^2Y^2(b-a)^2},\qquad K_b=\frac{X\cdot(x-w)\,(x-z)\cdot Y}{X^2Y^2\,ab},\notag\\
&K_c=\frac{X\cdot(x-w)\,Y\cdot(z-w)}{X^2Y^2Z^2\,a},\notag\\
&K_d=\frac{X\cdot Y\,(z-w)\cdot(x-z)\,a}{X^2Y^2Z^2(a-b)^2},\notag\\
&K_e=\frac{X\cdot Y\,(z-w)\cdot(x-w)\,b}{X^2Y^2Z^2(a-b)^2},\qquad K_f=\frac{X\cdot Y}{X^2Y^2Z^2},\notag\\
&K_{g1}=\frac{X\cdot(z-w)\,Y\cdot(x-w)}{X^2Y^2Z^2(a-b)},\qquad K_{g2}=\frac{X\cdot(x-w)\,Y\cdot(z-w)\,b}{X^2Y^2Z^2(a-b)a},\notag\\
&K_{g4}=\frac{X\cdot(z-w)\,Y\cdot(x-z)\,a}{X^2Y^2Z^2(a-b)b},\qquad K_{g5}=\frac{X\cdot(x-z)\,Y\cdot(z-w)}{X^2Y^2Z^2(a-b)},\notag\\
&\Sigma=K_{g1}-K_{g2}-K_{g4}+K_{g5}+K_e-K_d,\notag\\
&\KK_{\mathcal A}=\frac{1}{X^2Y^2Z^2}\Big[X\cdot(z-w)\,Y\cdot(x-w)+Y\cdot(z-w)\,X\cdot(x-z)\notag\\
&\hspace{6em}-\frac{c}{b}\,X\cdot(z-w)\,Y\cdot(x-z)-\frac{c}{a}\,Y\cdot(z-w)\,X\cdot(x-w)\Big],\notag\\
&\KK_4=\frac{1}{X^2Y^2Z^2}\Big[-\frac{X\cdot(z-w)\,Y\cdot(x-z)}{b}+\frac{Y\cdot(z-w)\,X\cdot(x-w)}{a}-X\cdot Y\Big],
\end{align}
with $c\equiv(x-w)\cdot(x-z)$. The longitudinal functions $\Lo^{\rm fin}$, $\Lop^{\rm fin}$, $\Lt$, $\Lf$ and the arctangent function $\mathcal A$ are defined in Table~\ref{tab:L}. The finite part of the sum of the three parts of $\mf B_3$ is
\begin{align}
\frac{d^3N^{\rm fin}_{\rm III}}{d^3k}={}&\pre\int_{x,y,z,x'}f^{abd}\,\frac{ig^4}{16\pi^5k^+}\Big[K_a\big(\Lo^{\rm fin}+\Lt-\lV+\lL\big)\notag\\
&\hspace{6em}-2K_b\,\Lop^{\rm fin}-2\lLb\,K_c\notag\\
&\hspace{6em}+2(\lL-\lV)(K_d-K_e)-2\lL\,K_f\notag\\
&\hspace{6em}-2\big(\Lo^{\rm fin}+\Lt\big)\Sigma+2\,\mathcal A\,\KK_{\mathcal A}+\Lf\,\KK_4\Big]\Ord_{\rm III}+{\rm c.c.},
\label{eq:finIII}
\end{align}
and the subtracted part is
\begin{align}
&\pre\int_{x,y,z,x'}f^{abd}\,\frac{ig^4}{16\pi^5k^+}\notag\\
&\qquad\times\Big\{\delta Y_T\big[K_b-2K_f\big]+\delta Y_P\big[-K_b+2K_c-2(K_{g1}-K_{g2}-K_{g4}+K_{g5})\big]\Big\}\Ord_{\rm III}+{\rm c.c.}
\label{eq:LLIII}
\end{align}
Three remarks.
\begin{itemize}
\item \emph{No power-like terms.} Each part of $\mf B_3$ separately produces terms proportional to $1/\Lambda$, $1/(k^+-\Lambda)$ and $1/(\vee-k^+)$ with the kernel $K_1=X\cdot Y/[X^2Y^2(b-a)]$. They come from the double pole $(p^+-k^+)^{-2}$ at the two edges of the excluded window. In units of $\frac{ig^4f^{abd}}{16\pi^5}K_1\Ord_{\rm III}$ they are
\begin{equation}
\underbrace{2\Big[\frac1\Lambda-\frac{1}{\vee-k^+}\Big]+2\Big[\frac1\Lambda-\frac{1}{k^+-\Lambda}\Big]}_{\text{instantaneous}}\ \underbrace{-2\Big[\frac1\Lambda-\frac{1}{\vee-k^+}\Big]}_{\text{non-inst., }p^+>k^+}\ \underbrace{-2\Big[\frac1\Lambda-\frac{1}{k^+-\Lambda}\Big]}_{\text{non-inst., }p^+<k^+}=0 .
\label{eq:powercancel}
\end{equation}
The cancellation is a direct consequence of the endpoint condition~\eqref{eq:endpoint}. With the opposite sign of the $p^+>k^+$ non-instantaneous term, the sum would be $\frac{4}{\Lambda}+\dots$.
\item \emph{Rapidity logarithms.} The logarithm $\log\frac{\vee-k^+}{\Lambda}$ that multiplies $K_a$ is compensated by the logarithms hidden in $\Lo$ and $\log\frac{\Lambda}{k^+-\Lambda}$, so $K_a$ appears in neither $\mathfrak T$ nor $\mathfrak P$. The term $K_b\,\delta Y_T$ comes from $p^+\to\Lambda$ in the commutator part of $\mf B_3$ and is a genuine region-T logarithm. The term $-2K_f\,\delta Y_T$ instead comes from the edge $p^+\to k^+-\Lambda$ of the excluded window. It is not a rapidity logarithm and must cancel against the diagonal row G1-I$'$.
\note{To be verified once G1-I$'$ is computed beyond LL.}
\item \emph{Limits.} For $\Lambda\to0$, $\vee\to\infty$: $\Lo^{\rm fin}+\Lt\to\log\frac{b}{a}$, $\Lop^{\rm fin}\to\log\frac ba$, $\Lf\to\log\frac{Z^2}{a}$, and the small logarithms~\eqref{eq:smalllogs} vanish. The finite part~\eqref{eq:finIII} converges to a constant; numerically, at a test configuration, the sequence $(\Lambda,\vee)=(10^{-3},10^2),\dots,(10^{-9},10^8)$ gives $-0.78910$, $-0.78204$, $-0.78197$, $-0.78197$ (App.~\ref{app:checks}).
\end{itemize}

\paragraph*{G1-IV ($\bar{\mf A}^\dagger\mf B_2$).} The colour operator is
\begin{align}
\Ord_{\rm IV}\equiv{}&\Big[U^{ab'}(x')\rho^{b'}(x')-U^{ab'}(w')\rho^{b'}(x')\Big]\notag\\
&\times\Big[f^{cdc'}U^{ac}(w)U^{bd}(z)U^{be}(y)\rho^e(y)\rho^{c'}(x)-f^{abc}U^{be}(y)U^{cd}(x)\rho^e(y)\rho^d(x)\Big].
\end{align}
The $p^+$ integral of the $\mf B_2$ structure appears again in G2-I and in G3-I, II. Define
\begin{align}
J_1^{ij}(x,z,w)&=\frac{(x-z)^j(z-w)^i}{Z^2\,b}-\frac{(x-w)^i(z-w)^j}{Z^2\,a}-\frac{(x-z)^j(x-w)^i}{ab}+\frac{\delta^{ij}}{2Z^2},\notag\\
J_2^{ij}(x,z,w)&=\KK^i(x-w)\Big[\KK^j(z-w)+\tfrac12\KK^j(x-z)\Big],\notag\\
\mathfrak P_{\mf B}^{ij}(x,z,w)&=\KK^j(x-z)\Big[\KK^i(z-w)-\tfrac12\KK^i(x-w)\Big].
\label{eq:J12}
\end{align}
The $p^+$ integral over $(\Lambda,\vee-k^+)$ gives $\frac{1}{k^+}\mathfrak F^{ij}$, with
\begin{equation}
\mathfrak F^{ij}=\Lp(a,b)\,J_1^{ij}-\frac12\log\frac{\vee}{\Lambda+k^+}\,\frac{\delta^{ij}}{Z^2}+\log\frac{\vee-k^+}{\Lambda}\,J_2^{ij}=\delta Y_T\,J_2^{ij}+\delta Y_P\,\mathfrak P_{\mf B}^{ij}+\mathfrak F^{{\rm fin},ij},
\label{eq:Fcal}
\end{equation}
\begin{equation}
\boxed{\ \mathfrak F^{{\rm fin},ij}=\Lp^{\rm fin}(a,b)\,J_1^{ij}+\frac12\lLb\,\frac{\delta^{ij}}{Z^2}+\lV\,J_2^{ij}.\ }
\label{eq:Ffin}
\end{equation}
$J_2$ is the soft gluon $z$ emitted by the measured gluon or by its source [region T, Eq.~\eqref{eq:B2soft}], and $\mathfrak P_{\mf B}$ is the measured gluon emitted by the harder gluon at $z$ [region P, Eq.~\eqref{eq:B2hard}]. The finite part of G1-IV is
\begin{equation}
\frac{d^3N^{\rm fin}_{\rm IV}}{d^3k}=\pre\int_{x,y,z,x'}\frac{ig^4}{8\pi^5k^+}\,\frac{(x'-w')^i(y-z)^j}{(x'-w')^2(y-z)^2}\,\mathfrak F^{{\rm fin},ij}(x,z,w)\,\Ord_{\rm IV}+{\rm c.c.},
\label{eq:finIV}
\end{equation}
and the subtracted part is the same with $\mathfrak F^{\rm fin}\to\delta Y_TJ_2+\delta Y_P\mathfrak P_{\mf B}$.

\paragraph*{G1-V ($\bar{\mf C}^\dagger\mf B_2$, the loop across the target).} Here the parent gluon $k^+$ splits into $p^+$ and $k^+-p^+$, so $0<p^+<k^+$ and there is no region P. The row has two parts.

\emph{Part 1} has the colour operator
\begin{align}
\Ord_{\rm V}\equiv{}&\Big[U^{ab'}(x')\rho^{b'}(x')-U^{ab'}(w')\rho^{b'}(x')\Big]\notag\\
&\times\Big[f^{adc}U^{db}(x)U^{ce}(z)\{\rho^e(y),\rho^b(y')\}-f^{abc}U^{cd}(y)U^{be}(y')\{\rho^d(y),\rho^e(y')\}\Big].
\end{align}
After the $w$ integration with the $\delta$-function of $\mf C_1$, the plus prescription makes the $p^+$ integrals finite:
\begin{align}
\frac{d^3N^{\rm fin}_{\rm V,1}}{d^3k}={}&\frac{1}{(2\pi)^3}\int_{x,z,y,x',y',w'}\frac{ig^4}{16\pi^4}\Bigg\{\frac{1}{k^{+2}}\int_0^{k^+}\frac{dp^+}{2\pi}e^{-ik\cdot(w'-z)}e^{-ik\cdot(z-x)\frac{p^+}{k^+}}\notag\\
&\qquad\times\frac{(x'-w')^i(z-x)^m(y-z)^l(y'-x)^j}{(x'-w')^2(z-x)^2(y-z)^2(y'-x)^2}\notag\\
&\qquad\times\Big[\delta_{im}\delta_{jl}-\Big[\frac{k^+}{p^+}\Big]_+\delta_{jm}\delta_{il}\Big]\notag\\
&-\frac{1}{k^+}\int_0^{k^+}\frac{dq^+}{2\pi}e^{-ik\cdot(w'-x)}e^{-ik\cdot(x-z)\frac{q^+}{k^+}}\notag\\
&\qquad\times\frac{(x'-w')\cdot(y'-x)\,(z-x)\cdot(y-z)}{(x'-w')^2(z-x)^2(y-z)^2(y'-x)^2}\Big[\frac{1}{q^+}\Big]_+\Bigg\}\Ord_{\rm V}\notag\\
&-\frac{1}{(2\pi)^3}\frac{ig^4}{32\pi^5k^+}\,\lL\int_{x,z,y,x',y',w'}\notag\\
&\qquad\times\Bigg[e^{-ik\cdot(w'-z)}\frac{(x'-w')\cdot(y-z)\,(z-x)\cdot(y'-x)}{(x'-w')^2(z-x)^2(y-z)^2(y'-x)^2}\notag\\
&\hspace{12em}+e^{-ik\cdot(w'-x)}\frac{(x'-w')\cdot(y'-x)\,(z-x)\cdot(y-z)}{(x'-w')^2(z-x)^2(y-z)^2(y'-x)^2}\Bigg]\Ord_{\rm V}+{\rm c.c.}
\label{eq:finV1}
\end{align}
The subtracted part is the last two lines with $\lL\to\delta Y_T$.

\emph{Part 2} has two colour structures. They arise from $f^{adc}f^{fbe}U^{db}(x)U^{ce}(z)$ by adding and subtracting $N_cU^{af}(x)$:
\begin{equation}
\Ord_{\rm V2a}\equiv\big[U^{ab'}(x')-U^{ab'}(w')\big]\big[-f^{adc}f^{fbe}U^{db}(x)U^{ce}(z)+N_cU^{af}(x)\big]\rho^{b'}(x')\rho^f(y),
\end{equation}
\begin{equation}
\Ord_{\rm V2b}\equiv\big[U^{ab'}(x')-U^{ab'}(w')\big]\,N_c\big[U^{af}(y)-U^{af}(x)\big]\rho^{b'}(x')\rho^f(y).
\end{equation}
With $\xi=p^+/k^+$ and $\mathcal D\equiv p^+(y-x)^2+(k^+-p^+)(y-z)^2$, define the six transverse brackets
\begin{align}
\mathcal B_1^{im}&=\frac{\delta^{im}d_\perp}{2(x-z)^2}\big[(y-x)^2-(y-z)^2\big],\notag\\
\mathcal B_2^{im}&=\frac{(y-x)^i(x-z)^m}{(x-z)^2}-\frac{(y-x)^i(y-z)^m}{2(y-z)^2},\notag\\
\mathcal B_3^{im}&=\delta_{im}\Big[\frac{(y-z)\cdot(x-z)}{(x-z)^2}+\frac{(y-x)\cdot(y-z)}{2(y-x)^2}-\frac{(y-x)^2-(y-z)^2}{2(x-z)^2}\Big],\notag\\
\mathcal B_4^{im}&=\frac{(y-z)^i(x-z)^m}{(x-z)^2}+\frac{(y-x)^m(y-z)^i}{2(y-x)^2},\notag\\
\mathcal B_5^{im}&=\delta_{im}\Big[\frac{(y-x)\cdot(x-z)}{(x-z)^2}-\frac{(y-x)\cdot(y-z)}{2(y-z)^2}-\frac{(y-x)^2-(y-z)^2}{2(x-z)^2}\Big],\notag\\
\mathcal B_6^{im}&=\frac{(y-x)^m(x-z)^i}{(x-z)^2}-\frac{(y-x)^m(y-z)^i}{2(y-z)^2}+\frac{(y-z)^m(x-z)^i}{(x-z)^2}+\frac{(y-x)^i(y-z)^m}{2(y-x)^2}.
\label{eq:B16}
\end{align}
The integrand of part 2 is then $\xi(1-\xi)\mathcal B_1-\frac{\xi}{1-\xi}\mathcal B_2+(1-\xi)\mathcal B_3-\frac{1-\xi}{\xi}\mathcal B_4+\xi\mathcal B_5-\mathcal B_6$, contracted with $\frac{(x'-w')^i(z-x)^m}{(x'-w')^2(z-x)^2}$ and divided by $\mathcal D$. Its finite part for the structure $\Ord_{\rm V2a}$ is
\begin{align}
\frac{d^3N^{\rm fin}_{\rm V2a}}{d^3k}={}&\frac{1}{(2\pi)^3}\frac{g^4}{8\pi^4k^+}\int_{x,y,z,x',w'}\frac{(x'-w')^i(z-x)^m}{(x'-w')^2(z-x)^2}\notag\\
&\times\Bigg\{\int_0^{k^+}\frac{dp^+}{2\pi}\frac{e^{-ik\cdot(w'-z)}e^{-ik\cdot(z-x)\xi}}{\mathcal D}\Big[\xi(1-\xi)\mathcal B_1+(1-\xi)\mathcal B_3-\Big[\frac1\xi\Big]_+\mathcal B_4+\mathcal B_4+\xi\mathcal B_5-\mathcal B_6\Big]\notag\\
&\quad+\int_0^{k^+}\frac{dq^+}{2\pi}\frac{e^{-ik\cdot(w'-x)}e^{-ik\cdot(x-z)\frac{q^+}{k^+}}}{(k^+-q^+)(y-x)^2+q^+(y-z)^2}\Big\{-\Big[\frac{k^+}{q^+}\Big]_+\mathcal B_2+\mathcal B_2\Big\}\Bigg\}^{im}\Ord_{\rm V2a}\notag\\
&-\frac{1}{(2\pi)^3}\frac{g^4}{16\pi^5k^+}\,\lL\int_{x,y,z,x',w'}\Bigg[e^{-ik\cdot(w'-z)}\frac{(x'-w')^i(z-x)^m\mathcal B_4^{im}}{(x'-w')^2(z-x)^2(y-z)^2}\notag\\
&\hspace{12em}+e^{-ik\cdot(w'-x)}\frac{(x'-w')^i(z-x)^m\mathcal B_2^{im}}{(x'-w')^2(z-x)^2(y-x)^2}\Bigg]\Ord_{\rm V2a}+{\rm c.c.},
\label{eq:finV2a}
\end{align}
and the subtracted part is the last two lines with $\lL\to\delta Y_T$. The structure $\Ord_{\rm V2b}$ has exactly the same kernel, so the same formula with $\Ord_{\rm V2a}\to\Ord_{\rm V2b}$ gives its finite part.

Unlike $\Ord_{\rm V2a}$, the operator $\Ord_{\rm V2b}$ does not vanish at $x=z$, and part~2b is UV divergent at $x\to z$. In Sec.~\ref{sec:rc} we show that the coefficient of this divergence is the integral of the gluon splitting function. This is the $\beta_0$ term of the loop across the target. Part~2b therefore has to be regularized in the same scheme as G1-II ($d_\perp=2-\epsilon$) before its finite part can be written as a number times LO plus a finite remainder.
\note{G1-V part~2b in dimensional regularization, and its overall weight (the factor 2 in Eq.~\eqref{eq:VV} together with the c.c.), remain to be completed.}
%TODO: complete G1-V(2b) in dim. reg.; fix its overall weight.

\paragraph*{G1-I$'$.} As discussed in Sec.~\ref{sec:groupI}, the diagonal region of $\bar{\mf B}_3^\dagger$ is known here only at leading log. Its finite part is not included.

%% -------------------------------------------------------------------
\subsection{Group II}\label{sec:finiteG2}

\paragraph*{G2-I, part 1 (the $\{\rho,\rho\}$ part of $\mathfrak B$ times the $f$ part).} The colour operators are
\begin{align}
\Ord^{(1)}&\equiv\Big[U^{ac'}(w')\{\rho^{c'}(x'),\rho^b(y)\}-U^{bc'}(z)U^{ad'}(x')U^{c'e'}(y)\{\rho^{d'}(x'),\rho^{e'}(y)\}\Big]\notag\\
&\qquad\times\Big[f^{cbd}U^{ac}(w)\rho^d(x)-f^{acd}U^{bc}(z)U^{de}(x)\rho^e(x)\Big],\notag\\
\Ord^{(2)}&\equiv\Big[f^{c'bd'}U^{ac'}(w')\rho^{d'}(x)-f^{ac'd'}U^{bc'}(z)U^{d'e'}(x)\rho^{e'}(x)\Big]\notag\\
&\qquad\times\Big[U^{ac}(w)\{\rho^c(x'),\rho^b(y)\}-U^{bc}(z)U^{ad}(x')U^{ce}(y)\{\rho^d(x'),\rho^e(y)\}\Big].
\end{align}
Both have the $\mf B_2$ structure of Eq.~\eqref{eq:Fcal}:
\begin{align}
\frac{d^3N^{\rm fin}_{\rm II.I,1}}{d^3k}={}&\pre\int_{x,y,z,x'}\Bigg\{-\frac{ig^4}{16\pi^5k^+}\frac{(x'-w')^i(y-z)^j}{(x'-w')^2(y-z)^2}\,\mathfrak F^{{\rm fin},ij}(x,z,w)\,\Ord^{(1)}\notag\\
&\hspace{8em}+\frac{ig^4}{16\pi^5k^+}\frac{(x'-w)^i(y-z)^j}{(x'-w)^2(y-z)^2}\,\mathfrak F^{{\rm fin},ij}(x,z,w')\,\Ord^{(2)}\Bigg\}.
\end{align}

\paragraph*{G2-I, part 2 (the $f$ part of $\mathfrak B$, squared).} This is the initial-state splitting term. Before the $p^+$ integration it reads
\begin{align}
\frac{d^3N_{\rm II.I,2}}{d^3k}={}&\frac{1}{(2\pi)^3}\frac{g^4}{4\pi^4}\int_{\Lambda}^{\vee-k^+}\frac{dp^+}{2\pi}\int_{z,w,w',x,x'}e^{-ik\cdot(w'-w)}\notag\\
&\times\frac{p^+k^+}{(p^++k^+)^2}\,\frac{\Phi^{ij}(x,z,w)\,\Phi^{ij}(x',z,w')}{\big[p^+b+k^+a\big]\big[p^+b'+k^+a'\big]}\,\Ord_{\rm II.I},
\label{eq:GIIIpart2}\\
\Phi^{ij}(x,z,w)\equiv{}&\frac{\delta^{ij}(b-a)}{2Z^2}+\frac{p^++k^+}{k^+}\Big[\frac{(x-z)^j(z-w)^i}{Z^2}-\frac{(x-z)^j(x-w)^i}{2a}\Big]\notag\\
&+\frac{p^++k^+}{p^+}\Big[\frac{(x-w)^i(z-w)^j}{Z^2}+\frac{(x-z)^j(x-w)^i}{2b}\Big],\notag
\end{align}
with the colour operator
\begin{align}
\Ord_{\rm II.I}\equiv{}&f^{c'bd'}f^{cbd}U^{ac'}(w')U^{ac}(w)\rho^{d'}(x')\rho^d(x)-f^{c'bd'}f^{acd}U^{ac'}(w')U^{bc}(z)U^{de}(x)\rho^{d'}(x')\rho^e(x)\notag\\
&-f^{ac'd'}f^{cbd}U^{ac}(w)U^{bc'}(z)U^{e'd'}(x')\rho^{e'}(x')\rho^d(x)+N_cU^{e'd}(x')U^{de}(x)\rho^{e'}(x')\rho^e(x).
\end{align}
$\Phi^{ij}$ is the transverse bracket of $\mf B_2$, Eq.~\eqref{eq:B2X}. Expanding the product $\Phi\Phi$ in powers of $(p^++k^+)$ gives five transverse kernels: $\KK_A$ multiplies $\frac{(p^++k^+)^2}{p^+k^+}$, $\KK_B$ multiplies $\frac{(p^++k^+)^2}{p^{+2}}$, $\KK_C$ multiplies $\frac{(p^++k^+)^2}{k^{+2}}$, $\KK_D$ multiplies $\frac{p^++k^+}{p^+}$, $\KK_E$ multiplies $\frac{p^++k^+}{k^+}$. There is also a term $\frac{d_\perp}{4Z^2Z'^2}(b-a)(b'-a')$:
\begin{align}
\KK_A={}&\frac{(x-w)\cdot(z-w')\,(x'-z)\cdot(z-w)}{Z^2Z'^2}-\frac{(x-w)\cdot(x'-w')\,(z-w)\cdot(x'-z)}{2Z^2a'}\notag\\
&+\frac{(x-z)\cdot(x'-z)\,(x-w)\cdot(z-w')}{2bZ'^2}\notag\\
&-\frac{(x'-z)\cdot(x-z)\,(x'-w')\cdot(x-w)}{4ba'}+\frac{(x'-w')\cdot(z-w)\,(x-z)\cdot(z-w')}{Z^2Z'^2}\notag\\
&+\frac{(x'-w')\cdot(z-w)\,(x-z)\cdot(x'-z)}{2Z^2b'}\notag\\
&-\frac{(x-z)\cdot(z-w')\,(x'-w')\cdot(x-w)}{2aZ'^2}-\frac{(x'-w')\cdot(x-w)\,(x'-z)\cdot(x-z)}{4ab'},\notag\\
\KK_B={}&\frac{(x'-w')\cdot(x-w)\,(z-w)\cdot(z-w')}{Z^2Z'^2}+\frac{(x'-w')\cdot(x-w)\,(x'-z)\cdot(z-w)}{2Z^2b'}\notag\\
&+\frac{(x-w)\cdot(x'-w')\,(x-z)\cdot(z-w')}{2bZ'^2}\notag\\
&+\frac{(x-z)\cdot(x'-z)\,(x'-w')\cdot(x-w)}{4bb'},\notag\\
\KK_C={}&\frac{(x'-z)\cdot(x-z)\,(z-w)\cdot(z-w')}{Z^2Z'^2}-\frac{(x'-z)\cdot(x-z)\,(x'-w')\cdot(z-w)}{2Z^2a'}\notag\\
&-\frac{(x-w)\cdot(z-w')\,(x-z)\cdot(x'-z)}{2aZ'^2}\notag\\
&+\frac{(x-z)\cdot(x'-z)\,(x'-w')\cdot(x-w)}{4aa'},\notag\\
\KK_D={}&\frac{(x-w)\cdot(z-w)\,(b'-a')}{2Z^2Z'^2}+\frac{(x-z)\cdot(x-w)\,(b'-a')}{4bZ'^2}\notag\\
&+\frac{(x'-w')\cdot(z-w')\,(b-a)}{2Z^2Z'^2}+\frac{(x'-z)\cdot(x'-w')\,(b-a)}{4Z^2b'},\notag\\
\KK_E={}&\frac{(x-z)\cdot(z-w)\,(b'-a')}{2Z^2Z'^2}-\frac{(x-z)\cdot(x-w)\,(b'-a')}{4aZ'^2}\notag\\
&+\frac{(x'-z)\cdot(z-w')\,(b-a)}{2Z^2Z'^2}-\frac{(x'-z)\cdot(x'-w')\,(b-a)}{4Z^2a'}.
\label{eq:KAE}
\end{align}
The $p^+$ integral is elementary. With $D\equiv a'b-ab'$ and the functions of Table~\ref{tab:L}, the finite part is
\begin{align}
\frac{d^3N^{\rm fin}_{\rm II.I,2}}{d^3k}={}&\frac{1}{(2\pi)^3}\frac{g^4}{8\pi^5k^+}\int_{z,w,w',x,x'}e^{-ik\cdot(w'-w)}\Bigg\{\frac{\KK_A}{D}\Big[\Lp^{\rm fin}(a,b)-\Lp^{\rm fin}(a',b')\Big]\notag\\
&+\KK_B\Big[\frac{\lV}{aa'}-\frac{b}{aD}\Lp^{\rm fin}(a,b)+\frac{b'}{a'D}\Lp^{\rm fin}(a',b')\Big]+\KK_C\Big[\frac{a'}{b'D}\Lp^{\rm fin}(a',b')-\frac{a}{bD}\Lp^{\rm fin}(a,b)\Big]\notag\\
&+\KK_D\Big[\frac{\lLb}{(b-a)(a'-b')}+\frac{b}{(b-a)D}\Lp^{\rm fin}(a,b)+\frac{b'}{(a'-b')D}\Lp^{\rm fin}(a',b')\Big]\notag\\
&+\KK_E\Big[-\frac{\lLb}{(b-a)(a'-b')}-\frac{a}{(b-a)D}\Lp^{\rm fin}(a,b)-\frac{a'}{(a'-b')D}\Lp^{\rm fin}(a',b')\Big]\notag\\
&+\frac{d_\perp}{4Z^2Z'^2}\Big[-\lLb\frac{aa'-bb'}{(b-a)(b'-a')}-\frac{ab(b'-a')}{(b-a)D}\Lp^{\rm fin}(a,b)+\frac{a'b'(b-a)}{(b'-a')D}\Lp^{\rm fin}(a',b')\notag\\
&\hspace{8em}+k^+\Big(\frac1\vee-\frac{1}{k^++\Lambda}\Big)\Big]\Bigg\}\Ord_{\rm II.I},
\label{eq:finGIIpart2}
\end{align}
and the subtracted part is
\begin{equation}
\frac{1}{(2\pi)^3}\frac{g^4}{8\pi^5k^+}\int_{z,w,w',x,x'}e^{-ik\cdot(w'-w)}\Big[\delta Y_T\,\frac{\KK_B}{aa'}+\delta Y_P\,\frac{\KK_C}{bb'}\Big]\Ord_{\rm II.I}.
\end{equation}
The coefficients of $\delta Y_P$ in the $\KK_A$, $\KK_B$, $\KK_D$, $\KK_E$ and $d_\perp$ terms vanish identically; only $\KK_C$ keeps one, since $\frac{a'}{b'D}-\frac{a}{bD}=\frac{1}{bb'}$. The kernel $\KK_B/(aa')$ is the soft gluon $z$ emitted by the measured gluon or its source (region T), and $\KK_C/(bb')$ is the measured gluon emitted by the harder gluon at $z$ (region P). The last term, $k^+(\frac1\vee-\frac{1}{k^++\Lambda})=-(1-\frac{k^+}{\vee})+\Ord(\Lambda)$, is a genuine $\Ord(1)$ term. Alone it is singular at $z\to w$ and $z\to w'$, and it must be kept together with the other $d_\perp$ terms, whose $\Lp$ have the opposite singularity. For $|z-w|,|z-w'|\ll|x-z|,|x'-z|$ the integrand~\eqref{eq:GIIIpart2} is proportional to the gluon splitting function, see Sec.~\ref{sec:dglap}.

\paragraph*{G2-II, G2-III.} Only four-charge terms. No finite part apart from the edge term of Sec.~\ref{sec:edge}.

\paragraph*{G2-IV (fragmentation of the scattered gluon).} With the colour operator
\begin{equation}
\Ord_{\rm II.IV}\equiv\big[U^{e'c}(y')-U^{e'c}(x')\big]\big[U^{ec}(y)-U^{ce}(x)\big]\rho^{e'}(x')\rho^e(x),
\end{equation}
the $\delta$-functions of $\mf C_1$ put the parent gluon at $y$ ($y'$) and the measured gluon at $w=y+\frac{p^+}{k^+}(y-z)$. Integrating over $w,w'$ gives
\begin{align}
\frac{d^3N_{\rm II.IV}}{d^3k}={}&\frac{1}{(2\pi)^3}\frac{g^4N_c}{4\pi^4k^+}\int_{\Lambda}^{\vee-k^+}\frac{dp^+}{2\pi}\int_{z,x,x',y,y'}e^{-ik\cdot(y'-y)\frac{k^++p^+}{k^+}}\,\frac{(y-z)^m(x-y)^l(y'-z)^{m'}(x'-y')^{l'}}{(y-z)^2(x-y)^2(y'-z)^2(x'-y')^2}\notag\\
&\times\Big[\frac{\delta_{lm}\delta_{l'm'}d_\perp p^+}{(k^++p^+)^2}+\frac{2}{k^+}\big(\delta_{ml'}\delta_{m'l}-\delta_{m'l'}\delta_{lm}\big)+\frac{p^+}{k^{+2}}\delta_{mm'}\delta_{ll'}+\frac{1}{p^+}\delta_{mm'}\delta_{ll'}\Big]\Ord_{\rm II.IV}.
\label{eq:GIIIV}
\end{align}
In terms of the momentum fraction of the measured gluon in the splitting, $\xi\equiv k^+/(k^++p^+)$, which runs from $k^+/\vee$ to $k^+/(k^++\Lambda)$, the transverse integral over $z$ produces a collinear pole. Its coefficient is the real gluon splitting function
\begin{equation}
\mathcal P_{gg}(\xi)\equiv2N_c\Big[\frac{\xi}{(1-\xi)_+}+\frac{1-\xi}{\xi}+\xi(1-\xi)\Big],
\label{eq:Pgg}
\end{equation}
convoluted with the LO spectrum at the parent momentum $k/\xi$. With $K_{yy'}\equiv\frac{(x-y)\cdot(x'-y')}{(x-y)^2(x'-y')^2}$, $\bar\epsilon$ the $\MSbar$ pole of the $z$ integral and $\mu$ its scale, the result is
\begin{align}
\frac{d^3N_{\rm II.IV}}{d^3k}={}&\frac{1}{(2\pi)^3}\frac{g^4N_c}{8\pi^4k^+}\,\delta Y_T\int_{x,x',y,y'}e^{-ik\cdot(y'-y)}\notag\\
&\qquad\times K_{yy'}\Big[\frac{1}{\bar\epsilon}-\log\big(4\pi^2\mu^2(y'-y)^2\big)\Big]\Ord_{\rm II.IV}\notag\\
&+\frac{1}{(2\pi)^3}\frac{g^4}{8\pi^4k^+}\int_{k^+/\vee}^{1}\frac{d\xi}{\xi^2}\int_{x,x',y,y'}e^{-ik\cdot(y'-y)/\xi}\notag\\
&\qquad\times K_{yy'}\,\frac{\mathcal P_{gg}(\xi)}{2}\Big[\frac{1}{\bar\epsilon}+1-\log\big(4\pi^2\mu^2(y'-y)^2\big)\Big]\Ord_{\rm II.IV}\notag\\
&-\frac{1}{(2\pi)^3}\frac{g^4N_c}{8\pi^4k^+}\int_{k^+/\vee}^{1}\frac{d\xi}{\xi^2}\int_{x,x',y,y'}e^{-ik\cdot(y'-y)/\xi}\notag\\
&\qquad\times\Big[\frac{(x-y)\cdot(y'-y)\,(x'-y')\cdot(y'-y)}{(x-y)^2(x'-y')^2(y'-y)^2}\,d_\perp\,\xi(1-\xi)\notag\\
&\hspace{14em}+K_{yy'}\Big(\frac{\xi}{(1-\xi)_+}+\frac{1-\xi}{\xi}\Big)\Big]\Ord_{\rm II.IV}.
\label{eq:GIIIVres}
\end{align}
Since $\Ord_{\rm II.IV}=-\Ord_{\rm LO}\big|_{w\to y,\,w'\to y'}$, the structure $K_{yy'}\Ord_{\rm II.IV}$ is the integrand of the LO spectrum~\eqref{eq:LOrelab} with the gluon at the position $y$ ($y'$) of the parent. The $1/\bar\epsilon$ term of the second line is the order-$g^2$ part of a convolution of the LO spectrum with the bare gluon fragmentation function $D_{g/g}(\xi)=\delta(1-\xi)-\frac{\alpha_s}{2\pi}\frac{1}{\bar\epsilon}\mathcal P_{gg}(\xi)+\Ord(\alpha_s^2)$. Absorbing it into the renormalized fragmentation function is the final-state DGLAP step (Sec.~\ref{sec:dglap}).
\note{The overall sign and power of $\xi$ in the identification with $\int d\xi\,D_{g/g}(\xi)\,d^3N_{\rm LO}(k/\xi)$ depend on whether $d^3N/d^3k$ or the boost-invariant $k^+d^3N/d^3k$ is convoluted, and on the orientation of the $\xi$ integral; to be fixed in the final version.}
%TODO: fix the normalization (xi power, orientation) of the fragmentation-function identification in G2-IV.

The decomposition~\eqref{eq:findef} of G2-IV is as follows. The first line is the region-T term, $\delta Y_T\,\mathfrak T_{\rm II.IV}$. A region-P term is hidden at the lower end $\xi\to k^+/\vee$ of the $\xi$ integrals: the $\frac{1-\xi}{\xi}$ term of $\mathcal P_{gg}$ combines with the Fourier transform of $\log(y'-y)^2$ at momentum $k/\xi$, which is $\propto\xi^2/\mathbf k^2$, into $\int d\xi/\xi=\delta Y_P+\dots$. Directly from Eq.~\eqref{eq:GIIIV}, before the $z$ integration, one finds for $p^+\gg k^+$
\begin{align}
\mathfrak P_{\rm II.IV}={}&\frac{1}{(2\pi)^3}\frac{g^4N_c}{8\pi^5k^+}\int_{z,w,w',x,x'}e^{-ik\cdot(w'-w)}\big[\KK(w-z)\cdot\KK(w'-z)\big]\notag\\
&\times\big[\KK(x-z)\cdot\KK(x'-z)\big]\notag\\
&\times\big[U^{e'c}(z)-U^{e'c}(x')\big]\big[U^{ec}(z)-U^{ce}(x)\big]\rho^{e'}(x')\rho^e(x).
\label{eq:PIIIV}
\end{align}
This is the region-P term in which the harder gluon at $z$ emits the soft measured gluon. The finite part of G2-IV is Eq.~\eqref{eq:GIIIVres} minus its first line minus $\delta Y_P\,\mathfrak P_{\rm II.IV}$.

%% -------------------------------------------------------------------
\subsection{Group III}\label{sec:finiteG3}

\paragraph*{G3-I, G3-II.} Both have the $\mf B_2$ structure with the prefactor $-\frac{ig^4}{8\pi^5k^+}$:
\begin{equation}
\frac{d^3N^{\rm fin}_{\rm III.I(II)}}{d^3k}=-\pre\int_{x,y,z,x'}\frac{ig^4}{8\pi^5k^+}\frac{(x'-w')^i(y-z)^j}{(x'-w')^2(y-z)^2}\,\mathfrak F^{{\rm fin},ij}(x,z,w)\,\Ord_{\rm III.I(II)}+{\rm c.c.},
\end{equation}
with
\begin{align}
\Ord_{\rm III.I}&\equiv\Big[U^{ab'}(x')U^{c'd'}(y)U^{c'b}(z)\rho^{b'}(x')\rho^{d'}(y)-U^{bd'}(z)U^{d'e'}(y)U^{ac'}(w')\rho^{c'}(x')\rho^{e'}(y)\Big]\notag\\
&\qquad\times\Big[f^{cbd}U^{ac}(w)\rho^d(x)-f^{acd}U^{bc}(z)U^{de}(x)\rho^e(x)\Big],\notag\\
\Ord_{\rm III.II}&\equiv\Big[U^{c'd'}(y)U^{ab'}(x')U^{bc'}(z)\rho^{d'}(y)\rho^{b'}(x')-U^{ab'}(x')\rho^b(y)\rho^{b'}(x')\Big]\notag\\
&\qquad\times\Big[f^{cbd}U^{ac}(w)\rho^d(x)-f^{acd}U^{bc}(z)U^{de}(x)\rho^e(x)\Big].
\end{align}
The subtracted parts are $-\frac{ig^4}{8\pi^5k^+}\frac{(x'-w')^i(y-z)^j}{(x'-w')^2(y-z)^2}[\delta Y_TJ_2^{ij}+\delta Y_P\mathfrak P_{\mf B}^{ij}]\Ord_{\rm III.I(II)}$.

\paragraph*{G3-III ($\mathfrak C^\dagger\mathfrak B$), part 1.} With the colour operators
\begin{align}
\Ord_4&\equiv U^{e'c'}(y')U^{d'b}(z)\rho^{e'}(x')-U^{bd'}(z)U^{c'e'}(x')\rho^{e'}(x'),\notag\\
\Ord_5&\equiv U^{ac}(w)\{\rho^c(y),\rho^b(x)\}-U^{bc}(z)U^{ad}(y)U^{ce}(x)\{\rho^d(y),\rho^e(x)\},
\end{align}
and the prefactor $f^{c'd'a}$, the $w'$ integration with the $\delta$-function of $\mf C_1$ gives $w'=y'+\frac{p^+}{k^+}(y'-z)$, and the plus prescription gives
\begin{align}
\frac{d^3N_{\rm III.III,1}}{d^3k}={}&\mathcal I_1+\log\frac{\vee-k^+}{\Lambda}\,\mathfrak T_1,\notag\\
\mathcal I_1\equiv{}&-\frac{1}{(2\pi)^3}\frac{ig^4f^{c'd'a}}{8\pi^4k^+}\int_0^{\vee-k^+}\frac{dp^+}{2\pi}\int_{w,x',y',z,y,x}e^{-ik\cdot(y'-w)}e^{-ik\cdot(y'-z)\frac{p^+}{k^+}}\notag\\
&\times\frac{(y'-z)^m(x'-y')^{l'}(y-w)^i(x-z)^j}{(y'-z)^2(x'-y')^2(y-w)^2(x-z)^2}\notag\\
&\times\Big[\frac{\delta_{l'm}\delta_{ij}}{p^++k^+}-\Big[\frac{1}{p^+}\Big]_+\delta_{jm}\delta_{il'}-\frac{1}{k^+}\delta_{im}\delta_{jl'}\Big]\Ord_4\Ord_5,\notag\\
\mathfrak T_1\equiv{}&\frac{1}{(2\pi)^3}\frac{ig^4f^{c'd'a}}{16\pi^5k^+}\int_{w,x',y',z,y,x}e^{-ik\cdot(y'-w)}\notag\\
&\times\frac{(y'-z)\cdot(x-z)\,(x'-y')\cdot(y-w)}{(y'-z)^2(x'-y')^2(y-w)^2(x-z)^2}\,\Ord_4\Ord_5.
\label{eq:GIIIIII1}
\end{align}
The region-P coefficient comes from the $-\frac{1}{k^+}\delta_{im}\delta_{jl'}$ term written in the original variable $w'$ (rule B):
\begin{equation}
\mathfrak P_1=\frac{1}{(2\pi)^3}\frac{ig^4f^{c'd'a}}{16\pi^5k^+}\int_{w,w',x',z,y,x}e^{-ik\cdot(w'-w)}\frac{(w'-z)\cdot(y-w)\,(x'-z)\cdot(x-z)}{(w'-z)^2(x'-z)^2(y-w)^2(x-z)^2}\,\Ord_4\big|_{y'\to z}\,\Ord_5.
\end{equation}
The finite part is
\begin{equation}
\frac{d^3N^{\rm fin}_{\rm III.III,1}}{d^3k}=\mathcal I_1+\big(\delta Y_P+\lV\big)\mathfrak T_1-\delta Y_P\,\mathfrak P_1+{\rm c.c.},
\label{eq:finGIIIIII1}
\end{equation}
and the subtracted part is $\delta Y_T\mathfrak T_1+\delta Y_P\mathfrak P_1$. The two explicit $\delta Y_P$ terms cancel the two logarithms hidden inside $\mathcal I_1$: $-\delta Y_P\mathfrak T_1$ from the plus prescription with its phase, Eq.~\eqref{eq:phase}, and $+\delta Y_P\mathfrak P_1$ from the $\frac{1}{k^+}$ term. Equivalently, the finite part is $\mathcal I_1$ with its large-$p^+$ tail subtracted, plus $\lV\mathfrak T_1$.

\paragraph*{G3-III, part 2.} With $\Ord_6\equiv f^{c'd'a}\Ord_4$ and $\Ord_7\equiv f^{cbd}U^{ac}(w)\rho^d(x)-f^{acd}U^{bc}(z)U^{de}(x)\rho^e(x)$, the same steps give $\frac{d^3N_{\rm III.III,2}}{d^3k}=\mathcal I_2+\log\frac{\vee-k^+}{\Lambda}\mathfrak T_2$, where
\begin{align}
\mathcal I_2\equiv{}&-\frac{1}{(2\pi)^3}\frac{g^4}{4\pi^4}\int_0^{\vee-k^+}\frac{dp^+}{2\pi}\int_{w,y',z,y,x}e^{-ik\cdot(y'-w)}e^{-ik\cdot(y'-z)\frac{p^+}{k^+}}\notag\\
&\times\frac{(y'-z)^m(x'-y')^{l'}}{(y'-z)^2(x'-y')^2}\,\frac{1}{p^+b+k^+a}\notag\\
&\times\Bigg[\frac{\delta_{l'm}d_\perp p^+(b-a)}{2Z^2(p^++k^+)^2}\notag\\
&\quad+\frac{p^+\delta_{l'm}}{k^+(p^++k^+)}\Big[\frac{(x-z)\cdot(z-w)}{Z^2}-\frac{(x-z)\cdot(x-w)}{2a}-\frac{b-a}{2Z^2}\Big]\notag\\
&\quad+\frac{\delta_{l'm}}{p^++k^+}\Big[\frac{(x-w)\cdot(z-w)}{Z^2}+\frac{(x-z)\cdot(x-w)}{2b}-\frac{b-a}{2Z^2}\Big]\notag\\
&\quad-\Big[\frac{1}{p^+}\Big]_+\Big[\frac{(x-w)^{l'}(z-w)^m}{Z^2}+\frac{(x-z)^m(x-w)^{l'}}{2b}\Big]\notag\\
&\quad-\frac{p^+}{k^{+2}}\Big[\frac{(x-z)^{l'}(z-w)^m}{Z^2}-\frac{(x-z)^{l'}(x-w)^m}{2a}\Big]\notag\\
&\quad-\frac{1}{k^+}\Big[\frac{(x-z)^m(z-w)^{l'}}{Z^2}-\frac{(x-z)^m(x-w)^{l'}}{2a}\notag\\
&\hspace{4em}+\frac{(x-w)^m(z-w)^{l'}}{Z^2}+\frac{(x-z)^{l'}(x-w)^m}{2b}\Big]\Bigg]\Ord_6\Ord_7,\notag\\
\mathfrak T_2\equiv{}&\frac{1}{(2\pi)^3}\frac{g^4}{8\pi^5k^+}\int_{w,y',z,y,x}e^{-ik\cdot(y'-w)}\notag\\
&\times\Big[\frac{(x-w)\cdot(x'-y')\,(z-w)\cdot(y'-z)}{Z^2a\,(y'-z)^2(x'-y')^2}+\frac{(x-z)\cdot(y'-z)\,(x-w)\cdot(x'-y')}{2ba\,(y'-z)^2(x'-y')^2}\Big]\Ord_6\Ord_7.
\label{eq:GIIIIII2}
\end{align}
The region-P coefficient comes from the term $-\frac{p^+}{k^{+2}}[\dots]$, which in the variable $w'$ behaves as $-\frac{1}{k^+p^+b}[\dots]$:
\begin{align}
\mathfrak P_2={}&\frac{1}{(2\pi)^3}\frac{g^4}{8\pi^5k^+}\int_{w,w',z,x,x'}e^{-ik\cdot(w'-w)}\big[\KK(x'-z)\cdot\KK(x-z)\big]\notag\\
&\times\KK^m(w'-z)\Big[\KK^m(z-w)-\tfrac12\KK^m(x-w)\Big]\Ord_6\big|_{y'\to z}\Ord_7 .
\end{align}
All other terms of the bracket fall faster than $1/p^+$. The finite part is
\begin{equation}
\frac{d^3N^{\rm fin}_{\rm III.III,2}}{d^3k}=\mathcal I_2+\big(\delta Y_P+\lV\big)\mathfrak T_2-\delta Y_P\,\mathfrak P_2+{\rm c.c.}
\end{equation}
This part contains the \emph{interference} of the initial- and final-state collinear splittings (Sec.~\ref{sec:dglap}).

\paragraph*{G3-IV.} Only four-charge terms. No finite part apart from the edge term.

\paragraph*{G3-V, G3-VI.} Both have exactly the longitudinal structure of G3-III part 1, with twice its prefactor ($-\frac{ig^4f}{4\pi^4k^+}$ instead of $-\frac{ig^4f}{8\pi^4k^+}$). For G3-V the prefactor colour factor is $f^{c'ca}$ and the colour operators are
\begin{align}
\Ord'_3\equiv{}& U^{e'c'}(y')\rho^{e'}(x')-U^{c'e'}(x')\rho^{e'}(x'),\notag\\
\Ord'_4\equiv{}& U^{ce}(x)U^{ad}(y)\rho^e(x)\rho^d(y)-U^{ad}(w)U^{ce}(x)\rho^e(x)\rho^d(y).
\end{align}
For G3-VI the colour factor is $f^{c'd'a}$ and
\begin{align}
\Ord''_1\equiv{}&U^{e'c'}(y')U^{d'b}(z)\rho^{e'}(x')-U^{bd'}(z)U^{c'e'}(x')\rho^{e'}(x'),\notag\\
\Ord''_5\equiv{}& U^{bc}(z)U^{ad}(y)U^{ce}(x)\rho^d(y)\rho^e(x)-U^{ad}(y)\rho^d(y)\rho^b(x).
\end{align}
With $\mathcal I_{\rm V}$, $\mathcal I_{\rm VI}$ the plus-prescription integrals,
\begin{equation}
\frac{d^3N^{\rm fin}_{\rm III.V(VI)}}{d^3k}=\mathcal I_{\rm V(VI)}+\big(\delta Y_P+\lV\big)\mathfrak T_{\rm V(VI)}-\delta Y_P\,\mathfrak P_{\rm V(VI)}+{\rm c.c.},
\end{equation}
where $\mathfrak T_{\rm V(VI)}$ and $\mathfrak P_{\rm V(VI)}$ are $\mathfrak T_1$ and $\mathfrak P_1$ with the prefactor doubled and the colour operators replaced by those of the row.

%% -------------------------------------------------------------------
\subsection{Edge terms and three-charge terms}\label{sec:edge}

\paragraph*{Four charges.} All four-charge terms, and the three-charge terms obtained by reordering them, multiply $\log(\vee/\Lambda)$ in the virtual rows, where $\Lambda<p^+<\vee$. In the real rows, where $\Lambda<p^+<\vee-k^+$, they multiply $\log\frac{\vee-k^+}{\Lambda}=\log\frac{\vee}{\Lambda}+\lV$. Since the four-charge terms cancel between real and virtual rows at LL (Sec.~\ref{sec:cancel}), the finite part contains the edge term
\begin{equation}
\frac{d^3N^{\rm fin}_{4\rho}}{d^3k}=\lV\times\big[\text{four-charge terms of the real rows}\big]=-\lV\times\big[\text{four-charge terms of the virtual rows}\big].
\end{equation}
It is $\Ord(k^+/\vee)$ and vanishes for $\vee\to\infty$; since $\vee$ is finite, we keep it.

\paragraph*{Three charges.} The finite parts of G1-III, G1-IV, G1-V(1), G2-I(1), G3-I, G3-II, G3-III(1), G3-V and G3-VI are products of three charges. After complete symmetrization each splits into a symmetric three-charge part and two-charge commutator terms. At LL the symmetric parts cancel. For the finite part there is no such requirement: symmetric three-charge terms are legitimate NLO contributions. They vanish for a Gaussian (McLerran--Venugopalan) projectile, in which the odd moments of $\rho$ vanish, but not, for example, for a projectile of a few valence quarks ($d^{abc}$ structures).

%% -------------------------------------------------------------------
\subsection{Summary of the finite part}\label{sec:finsum}

Table~\ref{tab:fin} lists where the finite part of the cross section resides.

\begin{table}[!tbp]
\centering
\renewcommand{\arraystretch}{1.2}
\begin{tabular}{L{3.7cm}L{3.5cm}L{7.4cm}}
\toprule
Row & finite part & content\\
\midrule
G1-I, G2-II, G2-III, G3-IV & edge term only & $\lV\times$(four-charge terms)\\
G1-II(1) & Eq.~\eqref{eq:finII1} & $\beta_0$ (running coupling), CMW constant, $\lV$ terms\\
G1-II(2) & Eq.~\eqref{eq:finII2} & transverse logarithms; no $\delta Y_P$ at all\\
G1-III & Eq.~\eqref{eq:finIII} & transverse logarithms, arctangent; no $1/\Lambda$\\
G1-IV, G2-I(1), G3-I, G3-II & $\mathfrak F^{\rm fin}$, Eq.~\eqref{eq:Ffin} & $\Lp^{\rm fin}\to\log\frac{(x-z)^2}{(x-w)^2}$, $\lV$\\
G1-V(1), V(2a) & Eqs.~\eqref{eq:finV1}, \eqref{eq:finV2a} & plus-prescription integrals over $0<p^+<k^+$\\
G1-V(2b) & Sec.~\ref{sec:rc} & UV divergent: $\beta_0$ from the loop across the target\\
G2-I(2) & Eq.~\eqref{eq:finGIIpart2} & initial-state collinear splitting (projectile DGLAP)\\
G2-IV & Eq.~\eqref{eq:GIIIVres} & fragmentation: $D_{g/g}$, $\mathcal P_{gg}$, minus the hidden $\delta Y_P\mathfrak P_{\rm II.IV}$\\
G3-III(1), V, VI & Eq.~\eqref{eq:finGIIIIII1} & plus-prescription integrals with explicit $\delta Y_P$ counterterms\\
G3-III(2) & Eq.~\eqref{eq:GIIIIII2} & interference of initial- and final-state splitting\\
G1-I$'$ & not computed & diagonal $\bar{\mf B}_3^\dagger$\\
\bottomrule
\end{tabular}
\caption{Where the finite part of the NLO cross section resides.}
\label{tab:fin}
\end{table}
